\section{Proof-producing schedules and simultaneous cut}
\label{sec:proof-schedule}

The trace compiler turns one bounded process into an object derivation.  A complete
mathematical proof contains many such components.  The proof schedule is the final
logical assembly layer.

\subsection{Schedule invariant}

Fix a permanent context \(\Gamma\) of input and boundary assumptions.  Let
\(D=[C_k,\ldots,C_1]\) be a reverse-topological list of contracts already compiled.
A schedule carries the invariant
\[
  \forall C\in D,\qquad \Gamma\vdash_{\PRAAx} C.
\]
To add a new node \(A\), provide a dependency sublist \(\Delta\subseteq D\) and a local
derivation
\[
  \Delta,\Gamma\vdash_{\PRAAx}A.
\]
Simultaneous cut replaces each dependency hypothesis by the corresponding stored
derivation, and each permanent hypothesis by itself.  The result is
\(\Gamma\vdash_{\PRAAx}A\).

\begin{theorem}[Finite proof schedule]
\label{thm:proof-schedule}
A finite topologically ordered list of local embedded derivations compiles to embedded
derivations of all node contracts from the permanent context.  In particular, the
root contract is derivable from exactly the listed boundary assumptions.
\end{theorem}

\begin{proof}
Induct on the node list.  The empty schedule is vacuous.  At a successor, apply the
simultaneous substitution/cut operation of the embedded deduction system to the local
proof.  The induction hypothesis supplies derivations for declared dependencies, and
ordinary hypotheses supply the permanent context.  Preserve all earlier derivations
unchanged.  The root is retrieved by membership.
\end{proof}

The \code{Schedule.extend} definition in the accompanying Lean source implements this
argument using \code{Prf.trans}.  It contains no semantic cast and no appeal to
soundness.

\subsection{Dependency declarations as an audit mechanism}

A local proof is accepted only under \(\Delta++\Gamma\), not under the entire global
collection of facts.  The inclusion proof \(\Delta\subseteq D\) is the acyclicity
certificate.  This has three benefits.

\begin{enumerate}[label=(\roman*),leftmargin=2.4em]
\item The generated proof term records the actual dependency graph.
\item A boundary assumption cannot leak into a supposedly unconditional component
unless it appears in \(\Gamma\) and the theorem statement.
\item Components can be replaced or strengthened locally, and unused earlier nodes do
not enter the final proof.
\end{enumerate}

For a generated library, one can additionally emit a machine-readable manifest mapping
each node to source files, theorem names, horizon transformers, and status.

\subsection{Proof-producing compilation versus proof checking}

A conventional certificate checker takes a numerical certificate and returns true or
false.  A proof-producing compiler instead constructs a syntax tree of deduction
rules.  The two may share the same finite algorithm, but their trust boundaries differ.

In the present architecture, the generated proof term can be large.  This is not a
logical defect.  Compression can be addressed later by factoring reusable lemmas,
using reflection theorems already proved in the object theory, or exporting to the
smaller native scalar calculus.  The first goal is a transparent proof with the right
axioms.

\subsection{Closed theorems and conditional theorems}

If \(\Gamma=[]\), the schedule yields a closed theorem.  If \(\Gamma\) contains a
modulus, gap, height bound, or other producer contract, the schedule yields a precise
conditional theorem.  When a producer is later proved, the schedules compose by one
more cut.  This makes partial progress cumulative: formalizing the finite verifier is
valuable even before the deep producer theorem is available.

\subsection{Native export}

After the ordinary derivation is complete, the repository theorem
\code{pra_iff_native} translates it to the native scalar core.  The final export stage
is therefore
\[
  \Gamma\vdash_{\PRAAx}\varphi
  \quad\Longrightarrow\quad
  \operatorname{ofPRA}(\Gamma)\vdash_{\mathrm{native}}
  \operatorname{ofPRA}(\varphi).
\]
This translation should occur only at the end.  The ordinary presentation has a much
richer checked arithmetic library, while the native core is the compact foundational
endpoint.

\subsection{A generated node template}

A practical compiler may emit one Lean declaration per node in the following pattern.

\begin{lstlisting}
-- Formula and fixed PR symbols for node v.
def nodeVFormula : Fm PTm n := ...

-- Local proof: only declared predecessors and base assumptions.
theorem nodeV_local : depsV ++ Gamma |-[(PRAAx)] nodeVFormula := by
  -- unfold fixed transition equations
  -- invoke Booleanization / bounded-search / arithmetic lemmas
  -- close a QF loop if needed
  ...

-- Add to proof schedule by simultaneous cut.
def scheduleV : Schedule Gamma (nodeVFormula :: done) :=
  Schedule.extend schedulePrev deps_subset nodeV_local
\end{lstlisting}

The notation in actual source follows the repository's scoped syntax.  The key property
is the type of \code{nodeV_local}, not the tactic script used to construct it.
